\section{Related Work}
One advantage of our framework is that it is highly related to existing work in the fields of linguistics and game theory. These fields are mostly concerned with asking how language develops rather than trying to engineer intelligent agents, however, it seems fruitful to leverage decades of work to inform our efforts. We give two notable examples next.

\paragraph{Game theory}
There is a large literature in game theory (\cite{farrell:1996}, \cite{blume:1998experimental}, \cite{crawford:1998survey}, \cite{Skyrms:2010}) and theoretical biology (\cite{nowak:2006evolutionary}) that focuses on signaling games. These are games of asymmetric information where one player ``knows'' something the other does not and has incentives to either reveal or not reveal it. One important strand of this literature concerns ``cheap talk'' games where agents must coordinate on some actions and have access to communication which initially has no meaning but gains meaning in the Nash equilibria of the game. Work in this literature focuses on understanding the theory of such games (\cite{farrell:1996}), whether learning or evolution can solve them (\cite{nowak:2006evolutionary}) and whether humans can learn to communicate and coordinate in laboratory settings (\cite{crawford:1998survey}, \cite{blume:1998experimental}).

\paragraph{Language evolution}
The field of language evolution is also closely related to the proposed framework. In the early and mid 1990's large groups of researchers were interested in multi-agent communication simulations using neural networks (\cite{Kirby:Hurford:2002},\cite{Steels:2015}) . While there is much there to be learned for our situation, given the focus of these works on studying the origins of language, the simulations from this period were rather restricted with respect to the nature and size of of input data, i.e, mostly artificial input and small networks (see ~\cite{Wagner:etal:2003} for a  comprehensive literature review on this field). In this position paper, we are arguing in favor of a multi-agent communication framework as the foundation of building machines that can communicate with us, thus focusing on more realistic data/input and less assumptions about their nature. 